% SUPERSEDED (2026-09-15): this file is NOT input by main.tex. Its four-bulk theorem and buffer bound were
% falsified by the audit (edges carry /psi_p instead of /d_lat; the count-times-timescale bound is a non-sequitur).
% The current theory lives in sec-theory.tex (main text) and sec-appendix-theory.tex (proofs). Kept only for history.
\section{Main-text design details}
\label{app:main-design-details}

\paragraph{Synthetic data.}
We generate Gaussian-mixture data whose signal lies in a
$\dint$-dimensional subspace embedded in a $\dlat$-dimensional
ambient space ($\dlat \geq \dint$), with $k$ cluster centers,
intra-cluster signal variance $\sigsig^2$ on the signal subspace,
and isotropic null-direction variance $\signoise^2$. Gaussian
mixtures give a closed-form population score (no noise floor),
an exact intrinsic/ambient split that lets us vary $\dlat$ and
$\dint$ independently, and the block-diagonal $\Sigmadata$ used by
the random-matrix-theory analysis.

\paragraph{Model families.}
The trainable score model is a 3-layer MLP with sinusoidal time
embedding, trained by score-matching MSE against the analytic
Gaussian-mixture score. Hidden width is held fixed across the
$\dlat$ sweep so parameter count and per-step FLOPs stay roughly
constant. The RFNN model follows \citet{bonnaire2025}: first-layer
features are frozen, only the readout is trained, and the resulting
quadratic score-matching loss admits a spectral analysis. We also
validate the picture empirically on MNIST and CelebA.

\paragraph{Timescales and sweeps.}
$\taugen$ is the first step at which test loss reaches within
$5\%$ of its per-run minimum. $\taumem$ is the first step at which
the \citet{somepalli2023diffusion} memorization fraction crosses
$1\%$, with runs censored at their training budget if the threshold
is not reached. For the RFNN at fixed $t$, generation is ill-defined,
so we use train-test divergence (gen-gap $>0.02$) as the
$\taumem$ proxy. We probe the buffer along two axes: varying
$\dlat$ at fixed $\dint=5$, and varying $\dint$ at fixed
$\dlat=20$.

\section{Experimental methods and design choices}
\label{app:methods}

\subsection{Data generation}
We generate anisotropic Gaussian-mixture data:
\begin{itemize}\itemsep1pt
\item Cluster centers $\{m_c\}_{c=1}^{k}$ are drawn iid uniformly on a sphere
      of radius $s$ inside the first $\dint$ coordinates of $\mathbb{R}^{\dlat}$;
      the remaining $\dlat - \dint$ ``null-space'' coordinates of every
      center are exactly zero.
\item For each training point $\mu$, a cluster index $c(\mu)$ is drawn
      uniformly, then $x^{\mu} = m_{c(\mu)} + Q\,\xi^{\mu}$ where $\xi$
      has independent components: variance $\sigsig^2 = 1$ along the
      $\dint$ signal dimensions and variance $\signoise^2$ along the
      $\dlat - \dint$ null-space dimensions. $Q$ is a fixed orthogonal
      basis whose first $\dint$ columns span the signal subspace.
\item Default values: $k = 10$, $\dint = 5$, $s = 3$,
      $\nsamp = 500$ training points, $2048$ held-out test points,
      $\sigsig = 1$, $\signoise \in \{0.01, 0.5\}$.
\end{itemize}

\subsection{Diffusion process}
We use the variance-preserving Ornstein-Uhlenbeck process with
$t \in [t_{\min}, t_{\max}] = [0.01, 3.0]$. Score targets are computed
analytically from the closed-form Gaussian-mixture density at noise level $t$;
this avoids the standard score-matching MSE loss noise floor and gives a
clean ``score error'' metric independent of training-set fluctuations.

\subsection{Models}
\paragraph{MLP.} 3-layer ReLU MLP with sinusoidal time embedding ($32$
frequency bands), input width $\dlat + 64$, output width $\dlat$.
Hidden width $h = 256$ held fixed across the $\dlat$ sweep. Adam with
learning rate $10^{-4}$, batch size $256$, $T = 3\times 10^5$ training
steps (extended to $5\times 10^6$ for the body timescale-scatter
figure), evaluation every $5{,}000$ steps with $5{,}000$ generated
samples.

\paragraph{RFNN.} Two-layer RFNN with $\tanh$ activation; first layer
$W \in \mathbb{R}^{\pwidth \times \dlat}$ is fixed at random Gaussian
initialization, the second-layer readout $A$ is trained. $\pwidth = 64\,\dlat$
to keep the rank-null bulk visible at every $\dlat$. Same optimizer, batch,
and evaluation cadence as the MLP.

\subsection{Evaluation metrics}
\begin{itemize}\itemsep1pt
\item \textbf{Score error}: integrated mean-squared error of the learned
      score against the analytic Gaussian-mixture score, averaged over
      $5{,}000$ test points.
\item \textbf{Memorization fraction}: fraction of generated samples whose
      nearest training-set neighbor is closer than one third of that
      neighbor's own nearest training-set
      neighbor~\citep{somepalli2023diffusion}.
\item \textbf{$\taugen$}: first step at which test loss is within $5\%$ of
      its per-run minimum.
\item \textbf{$\taumem$}: first step at which the memorization fraction
      exceeds $1\%$. Censored to ``$>$300k'' if not crossed within budget.
\end{itemize}

\paragraph{Why these definitions.}
We use the test-loss-plateau criterion for $\taugen$ rather than an FID-
or MMD-based criterion because (i) it is well-defined on both the MLP
and the RFNN, the latter trained at fixed $t$ with no usable generation
process, (ii) the analytic Gaussian-mixture score gives noise-floor-free
target values, removing the need for finite-sample FID estimates that
fluctuate with the number of generated samples, and (iii) the test-loss
plateau aligns directly with the moment when the population score has
been internalised. We use the \citet{somepalli2023diffusion}-style
memorization fraction for $\taumem$ on the MLP because it is the same
threshold \citet{bonnaire2025} use, allowing direct comparison with the
isotropic baseline. For the fixed-$t$ RFNN, no generation is available;
we substitute the gen-gap proxy ($\text{test loss} - \text{train loss}
> 0.02$), which monotonically tracks the canonical $\taumem$ on the MLP
and calibrates the two-track comparison.

\subsection{Per-experiment configurations}
\label{app:exp-configs}

Both sweeps run on the synthetic Gaussian-mixture data described
above. Shared
constants across every run: $\nsamp = 500$ training points, $2{,}048$
held-out test points, $k = 10$ cluster centers, $s = 3$, $\sigsig = 1$,
OU diffusion at $t \in [0.01, 3.0]$, seed $42$ for both data and the
random first layer $W$. Tables~\ref{tab:exp-configs}
and~\ref{tab:exp-hparams} list the sweeps and training
hyperparameters.

\begin{table}[!t]
\centering\small
\setlength{\tabcolsep}{4pt}
\caption{\textbf{Sweep configurations.} Both $\signoise$ values used
unless noted.}
\label{tab:exp-configs}
\begin{tabular}{@{}lllll@{}}
\toprule
Model & Vary & Range & Fixed & Width \\
\midrule
MLP                       & $\dlat$ & $\{5, 8, 10, 15, 20, 30, 40\}$ & $\dint = 5$ & $h = 256$ \\
RFNN                      & $\dlat$ & $\{5, 8, 10, 15, 20, 30, 40\}$ & $\dint = 5$ & $\pwidth = 64\,\dlat$ \\
RFNN (ext., $\signoise{=}0.5$) & $\dlat$ & $\{60, 80, 100, 150, 200\}$ & $\dint = 5$ & $\pwidth = 64\,\dlat$ \\
MLP                       & $\dint$ & $\{2, 5, 8, 12, 16, 20\}$ & $\dlat = 20$ & $h = 256$ \\
RFNN                      & $\dint$ & $\{2, 5, 8, 12, 16, 20\}$ & $\dlat = 20$ & $\pwidth = 1280$ \\
\bottomrule
\end{tabular}
\end{table}

\begin{table}[!t]
\centering\small
\setlength{\tabcolsep}{6pt}
\caption{\textbf{Training hyperparameters.}}
\label{tab:exp-hparams}
\begin{tabular}{@{}lll@{}}
\toprule
 & MLP & RFNN \\
\midrule
Optimizer & Adam & full-batch GD \\
Learning rate & $10^{-4}$ & closed-form \\
Batch size & $256$ & full ($\nsamp = 500$) \\
Total steps $T$ & $3 \times 10^5$ & $3 \times 10^5$ \\
Eval interval & $5{,}000$ & $5{,}000$ \\
Diffusion time & $t \sim \mathrm{Unif}[0.01, 3.0]$ & fixed $t = 0.01$ \\
Time embedding & sinusoidal, $32$ frequencies & n/a\\
Activation & ReLU & $\tanh$ \\
$\Umat$ Monte Carlo & n/a& $50$ noise samples per training point \\
\bottomrule
\end{tabular}
\end{table}

\subsection{Empirical bulk sizes}
\label{app:bulk-sizes}
Table~\ref{tab:bulk-sizes} reports the per-panel sizes of the four
detected bulks in Figure~\ref{fig:bulk-sizes-distances}, ordered
left-to-right by increasing eigenvalue (B1: rank-null tail; B4:
signal block). $\signoise = 0.01$, $\dint = 5$, $\nsamp = 500$,
$\pwidth = 64\,\dlat$.

\begin{table}[!t]
\centering\small
\setlength{\tabcolsep}{6pt}
\caption{\textbf{Empirical bulk sizes per $\dlat$.} B4 collapses to
$\dint = 5$ as predicted; B3 grows linearly with $\dlat - \dint$
(the noise-dim buffer); B1 absorbs the rank-null tail.}
\label{tab:bulk-sizes}
\begin{tabular}{@{}rrrrr@{}}
\toprule
$\dlat$ & B1 & B2 & B3 & B4 \\
\midrule
5  & 150  & 138 & 27 & 5 \\
8  & 344  & 133 & 30 & 5 \\
10 & 295  & 305 & 35 & 5 \\
15 & 699  & 216 & 40 & 5 \\
20 & 916  & 311 & 48 & 5 \\
30 & 1397 & 458 & 60 & 5 \\
40 & 1907 & 578 & 70 & 5 \\
\bottomrule
\end{tabular}
\end{table}

\subsection{Known pitfalls}
The $\tanh$ activation in the RFNN saturates whenever the data scale
divided by $\sqrt{\dlat}$ exceeds approximately $2$, collapsing all
eigenvalues into a degenerate single bulk. Configurations were chosen to
keep the operating point well inside the linear regime; see
Appendix~\ref{app:saturation}.

% ===== MLP block (sec 3 in the body) ===========================

\section{$\dlat$ sweep (synthetic data)}
\label{app:mlp-dlat-sweep}

The body sweeps $\dint$ at fixed $\dlat = 20$, the regime where
every configuration crosses the $1\%$ memorization threshold within
budget. The dual sweep varies $\dlat \in \{5, 8, 10, 15, 20, 30,
40\}$ at fixed $\dint = 5$ on the same Gaussian-mixture data,
hidden $= 256$, $\nsamp = 500$, with $\signoise = 0.5$ runs at 5M
steps and the $\signoise = 0.01$ counterparts at 300k. Many
high-ratio runs in this dual sweep do \emph{not} reach memorization
within budget, which is exactly the buffer mechanism prediction.

\begin{figure}[!t]
  \centering
  \includegraphics[width=0.85\linewidth]{figures/mlp_timescales_scatter.pdf}
  \caption{\textbf{$\dlat$ sweep: timescales versus $\dlat / \dint$
    at $\signoise = 0.5$.} 5M-step MLP. Triangles mark runs that do
    not reach $1\%$ memorization within budget. Dual of body
    Fig.~\ref{fig:mlp-timescale-gap}.}
  \label{fig:mlp-dlat-timescales}
\end{figure}

\begin{figure}[!t]
  \centering
  \includegraphics[width=0.7\linewidth]{figures/mlp_gen_gap_sn05.pdf}
  \caption{\textbf{$\dlat$ sweep: late generalization gap versus
    $\dlat / \dint$ at $\signoise = 0.5$.} Dual of body
    Fig.~\ref{fig:mlp-train-test}.}
  \label{fig:mlp-dlat-gen-gap}
\end{figure}

\begin{figure}[!t]
  \centering
  \includegraphics[width=0.7\linewidth]{figures/late_score_error_sn05.pdf}
  \caption{\textbf{$\dlat$ sweep: late score error versus
    $\dlat / \dint$ at $\signoise = 0.5$.} This is the wide-range
    score-error view moved out of the body.}
  \label{fig:mlp-dlat-late-score}
\end{figure}

\begin{figure}[!t]
  \centering
  \includegraphics[width=0.85\linewidth]{figures/mlp_score_error_curves_sn05.pdf}
  \caption{\textbf{$\dlat$ sweep: score error over training at
    $\signoise = 0.5$.} Per-step trajectories underlying the
    aggregate in Fig.~\ref{fig:mlp-dlat-late-score}, one curve per
    $\dlat$ (viridis colormap, low to high). Compare to
    Fig.~\ref{fig:mlp-dint-score-curves} (the $\dint$-sweep
    counterpart).}
  \label{fig:mlp-dlat-score-curves}
\end{figure}

\begin{figure}[!t]
  \centering
  \includegraphics[width=0.7\linewidth]{figures/mlp_timescales_bar_sn05.pdf}
  \caption{\textbf{$\dlat$ sweep: 300k-step bar-chart timescales at
    $\signoise = 0.5$.} The 300k-step bar-chart view of paired
    $\taugen$ (blue) and $\taumem$ (red), with hatched bars
    indicating $\taumem$ not reached within budget. The 5M-step
    extended-budget version is Fig.~\ref{fig:mlp-dlat-timescales}
    above.}
  \label{fig:mlp-dlat-timescales-bar}
\end{figure}

\section{$\dlat$ sweep at $\signoise = 0.01$}
\label{app:mlp-sn001}
\label{app:score-error-shape}

The $\signoise = 0.01$ counterparts of the $\dlat$ sweep figures
in Appendix~\ref{app:mlp-dlat-sweep}. The non-monotonic score-error
n-shape appears here: when the data/sample cliff is soft, increasing
$\dlat/\dint$ first worsens late score error and then partially
recovers at larger buffer widths. The peak is visible in
Fig.~\ref{fig:mlp-dlat-late-score-sn001}.

\begin{figure}[!t]
  \centering
  \includegraphics[width=0.7\linewidth]{figures/mlp_timescales_bar_sn001.pdf}
  \caption{\textbf{$\signoise = 0.01$ counterpart of
    Fig.~\ref{fig:mlp-dlat-timescales-bar}.}}
  \label{fig:mlp-dlat-timescales-sn001}
\end{figure}

\begin{figure}[!t]
  \centering
  \includegraphics[width=0.7\linewidth]{figures/mlp_gen_gap_sn001.pdf}
  \caption{\textbf{$\signoise = 0.01$ counterpart of
    Fig.~\ref{fig:mlp-dlat-gen-gap}.}}
  \label{fig:mlp-dlat-gen-gap-sn001}
\end{figure}

\begin{figure}[!t]
  \centering
  \includegraphics[width=0.7\linewidth]{figures/mlp_score_error_curves_sn001.pdf}
  \caption{\textbf{$\signoise = 0.01$ counterpart of
    Fig.~\ref{fig:mlp-dlat-score-curves}.}}
  \label{fig:mlp-dlat-score-curves-sn001}
\end{figure}

\begin{figure}[!t]
  \centering
  \includegraphics[width=0.7\linewidth]{figures/late_score_error_sn001.pdf}
  \caption{\textbf{$\signoise = 0.01$ counterpart of
    Fig.~\ref{fig:mlp-dlat-late-score}.} The non-monotonic peak
    discussed in Sec.~\ref{sec:discussion} is visible here.}
  \label{fig:mlp-dlat-late-score-sn001}
\end{figure}

\section{$\dint$ sweep additional figures}
\label{app:mlp-dint}

The body $\dint$-sweep figures (Figs.~\ref{fig:mlp-timescale-gap},
\ref{fig:mlp-train-test}) report aggregate timescales and gen-gap;
this section collects the per-step score-error trajectories and the
$\dint$-sweep late score error, complementing the wider-range
$\dlat$-sweep view in Fig.~\ref{fig:mlp-dlat-late-score}.

\begin{figure}[!t]
  \centering
  \includegraphics[width=0.85\linewidth]{figures/mlp_dint_score_curves_sn05.pdf}
  \caption{\textbf{$\dint$ sweep: score error over training at
    $\signoise = 0.5$.} One curve per $\dlat / \dint$ ratio
    (viridis colormap). Per-step trajectories underlying the
    body's $\dint$-sweep aggregates.}
  \label{fig:mlp-dint-score-curves}
\end{figure}

\begin{figure}[!t]
  \centering
  \includegraphics[width=0.7\linewidth]{figures/mlp_dint_late_score_sn05.pdf}
  \caption{\textbf{$\dint$ sweep: late score error versus
    $\dlat / \dint$ at $\signoise = 0.5$.} $\dint$-sweep counterpart
    of Fig.~\ref{fig:mlp-dlat-late-score}.}
  \label{fig:mlp-dint-late-score}
\end{figure}

\section{$\dint$ sweep at $\signoise = 0.01$}
\label{app:mlp-dint-sn001}

The $\signoise = 0.01$ counterparts of the body's $\dint$ sweep at
fixed $\dlat = 20$. These runs are at 300k training steps (the 5M
extension was only run at $\signoise = 0.5$).

\begin{figure}[!t]
  \centering
  \includegraphics[width=0.7\linewidth]{figures/mlp_dint_timescales_sn001.pdf}
  \caption{\textbf{$\signoise = 0.01$ counterpart of body
    Fig.~\ref{fig:mlp-timescale-gap}.}}
  \label{fig:mlp-dint-timescales-sn001}
\end{figure}

\begin{figure}[!t]
  \centering
  \includegraphics[width=0.7\linewidth]{figures/mlp_dint_gen_gap_sn001.pdf}
  \caption{\textbf{$\signoise = 0.01$ counterpart of body
    Fig.~\ref{fig:mlp-train-test}.}}
  \label{fig:mlp-dint-gen-gap-sn001}
\end{figure}

\begin{figure}[!t]
  \centering
  \includegraphics[width=0.85\linewidth]{figures/mlp_dint_score_curves_sn001.pdf}
  \caption{\textbf{$\signoise = 0.01$ counterpart of
    Fig.~\ref{fig:mlp-dint-score-curves}.}}
  \label{fig:mlp-dint-score-curves-sn001}
\end{figure}

\begin{figure}[!t]
  \centering
  \includegraphics[width=0.7\linewidth]{figures/mlp_dint_late_score_sn001.pdf}
  \caption{\textbf{$\signoise = 0.01$ counterpart of
    Fig.~\ref{fig:mlp-dint-late-score}.}}
  \label{fig:mlp-dint-late-score-sn001}
\end{figure}

\section{Real-data experiments}
\label{app:real-data}

\subsection{Setup}
We test the buffer mechanism on two real-image datasets, MNIST and
CelebA, with the latent space learned end-to-end by a VAE.

\paragraph{Intrinsic dimension.} MLE-based estimates on raw pixels
\citep{pope2021} place MNIST at $\dint \approx 13$ and CelebA at
$\dint \approx 26$. Our $\dlat$ sweeps---$\{10, 15, 20, 25, 30, 40\}$
for MNIST and $\{20, 25, 30, 35, 40, 50\}$ for CelebA---bracket these
values to probe the $\dlat > \dint$ buffer regime predicted by the
synthetic theory. Alternative estimators
\citep{facco2017twonn, stanczuk2024manifold, kamkari2024geometric}
disagree at the $2{\times}$ level, so the specific numerical
predictions of buffer width $\dlat - \dint$ are
estimator-conditional (Sec.~\ref{sec:discussion}).

\paragraph{VAE.} Trained on half the dataset (the ``encoder half''),
varying the bottleneck $\dlat$ across the same sweep used in
synthetic experiments. The remaining half (the ``diffusion half'')
is reserved for diffusion training and evaluation.

\paragraph{Diffusion training set.} Each diffusion model is trained
on only $\nsamp = 1000$ latents drawn from the diffusion half, kept
small to put the model in the memorization-prone overparameterized
regime.

\paragraph{Score model.} Identical architecture to the synthetic
MLP (Sec.~\ref{app:methods}): 3-layer ReLU MLP with sinusoidal time
embedding, hidden width $h = 256$ held fixed across the $\dlat$
sweep, same optimizer, batch size, and step budget.

\paragraph{Metrics.}
\begin{itemize}\itemsep1pt
\item \textbf{$\taugen$}: first step at which FID is within $5\%$
      of its per-run minimum, i.e.\ the FID plateau is reached.
      We use FID here rather than the synthetic test-loss-plateau
      criterion because for real data FID is the canonical
      sample-quality signal.
\item \textbf{$\taumem$}: first step at which the gen-gap
      proxy~($\text{test loss} - \text{train loss} > 0.02$) is
      crossed; we use the gen-gap proxy here rather than the
      \citet{somepalli2023diffusion}-style memorization fraction
      because the pixel-space nearest-neighbor threshold rarely
      triggers at $\nsamp = 1000$ in the latent space and the
      gen-gap proxy monotonically tracks the canonical $\taumem$
      on synthetic data.
\end{itemize}

\subsection{Results}
The three synthetic findings transfer to real data at
$\nsamp = 1000$. Fig.~\ref{fig:real-timescale-gap} shows that
$\taumem$ grows with $\dlat$ on both datasets, going from $10\,k$
to $80\,k$ steps on MNIST and from $30\,k$ to $440\,k$ steps on
CelebA across the sweep---a $\sim$$10{\times}$ delay in
memorization onset attributable to the noise-dim buffer.
Fig.~\ref{fig:real-train-test} reports the late gen-gap.
Fig.~\ref{fig:real-min-fid} shows minimum FID is roughly flat
across the sweep, so the delay benefit of larger $\dlat$ does not
come at the cost of generation quality;
Fig.~\ref{fig:real-fid} reports the per-step FID trajectories.

\begin{figure}[!t]
  \centering
  \includegraphics[width=\linewidth]{figures/real_timescales.pdf}
  \caption{\textbf{Real data: timescales versus $\dlat / \dint$.}
    5M-step MLP, $\nsamp = 1000$. $\taugen$ (blue) is the first
    step at which test loss is within $5\%$ of its per-run minimum;
    $\taumem$ (red) is the first step at which the gen-gap proxy
    crosses $0.02$. $\taumem$ grows by roughly $10{\times}$ across
    each sweep ($10\,k \to 80\,k$ on MNIST, $30\,k \to 440\,k$ on
    CelebA), tracking the buffer-mechanism prediction.}
  \label{fig:real-timescale-gap}
\end{figure}

\begin{figure}[!t]
  \centering
  \includegraphics[width=\linewidth]{figures/real_late_gen_gap.pdf}
  \caption{\textbf{Real data: late generalization gap versus
    $\dlat / \dint$.} Mean of $\text{test loss} - \text{train loss}$
    over the last $5\%$ of training.}
  \label{fig:real-train-test}
\end{figure}

\begin{figure}[!t]
  \centering
  \includegraphics[width=\linewidth]{figures/real_min_fid.pdf}
  \caption{\textbf{Real data: minimum FID versus $\dlat / \dint$.}
    Sample quality is roughly preserved across the sweep, so the
    delay-in-memorization benefit of larger $\dlat$ does not come
    at the cost of generation quality.}
  \label{fig:real-min-fid}
\end{figure}

\begin{figure}[!t]
  \centering
  \includegraphics[width=\linewidth]{figures/real_fid_curves.pdf}
  \caption{\textbf{Real data: FID over training.} Per-step
    trajectories underlying Fig.~\ref{fig:real-min-fid}, one curve
    per $\dlat / \dint$ ratio.}
  \label{fig:real-fid}
\end{figure}

% ===== RFNN block (sec 4 in the body) ==========================

\section{RFNN: sketch of an analytical derivation}
\label{app:proof-sketch}

\subsection{Four-bulk intuition}
\label{app:rfnn-fourbulk-intuition}

\citet{bonnaire2025} note that the population bulk $\rho_2$
develops internal structure when $\Sigmadata$ has multiple
eigenvalues, but at signal-to-null ratio $3{\times}$ they treat it
as a single deformed bulk. The latent-diffusion regime sits at the
opposite extreme ($\sigsig^2/\signoise^2 \approx 10^4$ at
$\signoise = 0.01$, $\approx 4$ at $\signoise = 0.5$): the humps
inside $\rho_2$ separate into well-isolated bulks decades apart,
which is why we count them as distinct.

\paragraph{What the gen-gap measures.}
The training set carries two kinds of structure. The data bulks
correspond to \emph{population} structure: directions of variance
that exist in the population and happen to also be visible in the
training sample. The sample bulk corresponds to \emph{null
structure}: directions that look like signal in the $\nsamp$
training points but do not exist in the population. The RFNN
absorbs eigenmodes in order of decreasing $\lambda_i$, so it fits
the population structure first and the null structure later. While
it is fitting population structure, train and test loss drop
together. Once it begins fitting null structure, train loss keeps
dropping but test loss does not, and the gap opens. The gen-gap is
therefore a direct measure of how much null structure the readout
has absorbed, which is why it opens at $\taumem$ and why the
noise-dim buffer also delays its onset: widening the buffer pushes
back the moment when null structure begins to be learned.

\paragraph{Decomposing the four-bulk structure.}
The split is easiest to see in sorted-index rather than density
language. In the isotropic setting, all population directions have
the same scale, so the data-supported part of the spectrum appears
as one bulk. In our block-anisotropic setting, the signal coordinates
and the null coordinates have different variances before the random
feature lift. The Hermite-linear part of the $\tanh$ features
therefore produces two data-side eigenvalue scales: a high-eigenvalue
signal block of width $\dint$ and a lower-eigenvalue null-direction
block of width $\dlat-\dint$.

Equivalently, the four populations can be viewed along two axes:
shared across the training set versus sample-specific, and signal
subspace versus null subspace. The four resulting groups are:
(i)~\textsc{bulk\_data\_signal}, the $\dint$ meaningful signal
directions shared by the population;
(ii)~\textsc{bulk\_data\_noise}, the $\dlat-\dint$ null directions
that carry small but population-level variance;
(iii)~\textsc{bulk\_sample\_signal}, sample-specific fluctuations
inside the signal subspace; and
(iv)~\textsc{bulk\_sample\_noise}, sample-specific fluctuations
inside the null subspace. The first two refine Bonnaire's
population bulk $\rho_2$, while the latter two make up the
sample-specific bulk $\rho_1$.

The signal and noise-dim bulks of $\Umat$ inherit the block
structure of $\Sigmadata$ directly: $X^\top X / \nsamp$ on the same
training data has $\dint$ large eigenvalues in the signal subspace
and $\dlat - \dint$ small eigenvalues in the null directions,
separated by roughly $\sigsig^2 / \signoise^2$, and the lift by $W$
carries them into the data bulks of $\Umat$ up to a constant from
$W$'s spectrum. The sample and rank-null bulks are feature-space
artifacts of $\nsamp$ training points lifted into
$\pwidth > \dlat$ random features. The buffer mechanism is thus
gradient flow following the data's own variance hierarchy. As
$\dlat$ grows at fixed $\dint$, the signal block keeps the same
width, but the null-direction block widens. In the sorted spectrum,
the cliff associated with sample-specific fitting is pushed to later
indices, while the signal-learning cliff stays tied to the intrinsic
dimension. Thus the extra latent coordinates form a trainable
population buffer: they absorb optimization time after the model has
learned signal directions but before it reaches sample-specific
directions.

\subsection{Main-text theoretical statement}
\label{app:rfnn-theoretical-statement}

Using a Hermite expansion of the RFNN feature map, the four-bulk
picture predicts that in the proportional limit
$\dlat, \nsamp, \pwidth \to \infty$, the spectrum of $\Umat$
separates into four contiguous bulks with counts
\[
  \dint,\qquad \dlat-\dint,\qquad \nsamp,\qquad \pwidth-\dlat-\nsamp,
\]
corresponding respectively to data-signal, data-noise, sample, and
rank-null modes. Combining this spectral decomposition with the
per-mode decay in Eq.~\eqref{eq:rfnn-mode-decay} gives the buffer
prediction: increasing $\dlat-\dint$ adds modes between population
learning and sample-specific memorization. Thus $\taumem$ is delayed
while $\taugen$ remains approximately fixed.

Equivalently, the memorization delay scales with the width of the
noise-dimension bulk,
\begin{equation}
\begin{aligned}
  \taumem - \taugen
  &\;\gtrsim\; (\dlat - \dint)\,
    \frac{1}{\lambda_{\min}^{\mathrm{noise\text{-}dim}}} \\
  &\;\approx\; (\dlat - \dint)\,
    \frac{\psi_p}{\mu_1^{\,2}\beta_t^{\,2}} .
\end{aligned}
\label{eq:buffer-bound}
\end{equation}
This is the formal version of the noise-dimension buffer mechanism:
larger excess latent dimensionality lengthens the path from
generalization to memorization without changing the leading signal
timescale.

\subsection{Derivation sketch}

We extend \citeauthor{bonnaire2025}'s replica computation to block-diagonal data covariance
\(\Sigmadata = \diag\!\bigl(\sigsig^2\mathbb{I}_{\dint},\,
                           \signoise^2\mathbb{I}_{\dlat - \dint}\bigr)\).
Throughout, $x^\mu = m_{c(\mu)} + \xi^\mu$ with $m_{c(\mu)}$ a cluster
center in the $\dint$-dimensional signal subspace and
$\xi^\mu \sim \mathcal{N}(0, \Sigmadata)$. The diffused state at time
$t$ is $x_t^\mu = e^{-t}\,x^\mu + \sqrt{\Delta_t}\,\eta^\mu$ with
$\Delta_t = 1 - e^{-2t}$ and $\eta^\mu \sim \mathcal{N}(0, \mathbb{I}_{\dlat})$.
We work in the proportional limit $\dlat,\nsamp,\pwidth\to\infty$ with
$\psi_p\!\equiv\!\pwidth/\dlat$ and $\psi_n\!\equiv\!\nsamp/\dlat$ fixed
and $\psi_p > 1 + \psi_n$ (so the rank-null tail has room to exist).

\paragraph{Hermite expansion of $\tanh$.}
For data scale $s$ small compared with $\sqrt{\dlat}$
(Appendix~\ref{app:saturation}), the pre-activations
$Wx_t^\mu/\sqrt{\dlat}$ have order-one variance per coordinate and stay
inside the analytic regime of $\tanh$. Let
$\mu_k = \E_{z\sim\mathcal{N}(0,1)}[h_k(z)\tanh(z)]$ denote the
Hermite coefficients of $\tanh$ against the probabilists' Hermite
polynomials $\{h_k\}$; $\tanh$ is odd so $\mu_k = 0$ for even $k$,
and numerically $\mu_1 \approx 0.606$, $\mu_3 \approx -0.099$.
For each pair of inputs $x,y\in\R^{\dlat}$ the random vector
$((Wx)_a,(Wy)_a)$ is jointly Gaussian with mean zero and covariance
$\diag(\|x\|^2,\|y\|^2)/\dlat$ plus off-diagonal $x^\top y/\dlat$,
\emph{regardless of whether $x$ or $y$ has nonzero mean as a sample}
(the cluster centers shift the distribution of $x$, not the
conditional Gaussian over $W$ at fixed $x$). Consequently the Hermite
expansion of $\tanh$ against this Gaussian gives $\mu_0=0$
(odd function against a zero-mean Gaussian), and the pairwise kernel
admits the proportional-limit decomposition
\citep{pennington2017nonlinear, benigni2021eigenvalue}
\begin{equation}
  \E_W\!\bigl[\tanh\!\bigl((Wx)_a\bigr)\tanh\!\bigl((Wy)_a\bigr)\bigr]
   \;=\; \mu_1^{\,2}\,\frac{x^\top y}{\dlat}
        \;+\; \eta_\star(x,y)\,\mathbb{1}\!\bigl[x=y\bigr]
        \;+\; o_d(1),
  \label{eq:kernel-decomp}
\end{equation}
where $\eta_\star(x,x) = \sum_{k\ge 3}\mu_k^{\,2}\,(\|x\|^2/\dlat)^k$
is the higher-Hermite mass that survives only on the diagonal.
\citet{ELKaroui2010spectrum} establishes the same off-diagonal/diagonal
split at the level of kernel random matrices.



\paragraph{Two-piece decomposition of $\Umat$.}
Substituting \eqref{eq:kernel-decomp} into
$\Umat = \tfrac{1}{\nsamp}\sum_\mu \E_\eta[\phi(x_t^\mu)\phi(x_t^\mu)^\top]$,
$\phi(x) = \tfrac{1}{\sqrt{\pwidth}}\tanh(Wx/\sqrt{\dlat})$, and
averaging over the diffusion noise yields, to leading order in $1/\dlat$,
\begin{equation}
  \Umat \;=\; \Umat^{\mathrm{lin}} + \Umat^{\mathrm{diag}} + R_t,
  \qquad \|R_t\|_{\mathrm{op}} = o_d(1),
  \label{eq:U-split}
\end{equation}
where
\begin{align}
  \Umat^{\mathrm{lin}}
    \;&=\; \frac{\mu_1^{\,2}}{\pwidth\,\dlat}\,W\,M_t\,W^\top,
    \label{eq:U-lin}\\
  \Umat^{\mathrm{diag}}
    \;&=\; \frac{1}{\pwidth\,\nsamp}\sum_{\mu=1}^{\nsamp}
           \eta_\star\!\bigl(x_t^\mu,\,x_t^\mu\bigr)\,
           \phi^{\perp}\!(x_t^\mu)\,\phi^{\perp}\!(x_t^\mu)^\top,
    \label{eq:U-diag}
\end{align}
and $\phi^\perp(x)$ is the residual feature vector after projecting
out the Hermite-1 component. By construction
$\E_W[\phi^\perp(x)\phi^\perp(y)^\top] = 0$ for $x\neq y$ and
$\E_W[\|\phi^\perp(x)\|^2] = \eta_\star(x,x)+o_d(1)$, so the rank-one
outer products in \eqref{eq:U-diag} are mutually $W$-orthogonal in
the proportional limit.

\paragraph{Block decomposition of $M_t$.}
Direct computation gives
\begin{equation}
  M_t \;=\; \tfrac{1}{\nsamp}\sum_\mu \E_\eta\!\bigl[x_t^\mu(x_t^\mu)^\top\bigr]
        \;=\; e^{-2t}\bigl(\hat C + \hat S\bigr) + \Delta_t\,\mathbb{I}_{\dlat},
  \label{eq:Mt}
\end{equation}
with $\hat C = \tfrac{1}{\nsamp}\sum_\mu m_{c(\mu)}m_{c(\mu)}^\top$ the
empirical centroid Gram and $\hat S = \tfrac{1}{\nsamp}\sum_\mu
\xi^\mu(\xi^\mu)^\top$ the empirical within-cluster covariance.
In population, $\E\hat C$ has rank $\dint$ with eigenvalues
$\Theta(s^2/\dint)$ in the signal subspace and $\E\hat S = \Sigmadata$
contributes $\sigsig^2$ on the signal block and $\signoise^2$ on the
null block. Writing $\alpha_t^2 \equiv e^{-2t}(s^2/\dint + \sigsig^2) + \Delta_t$
for the signal-block eigenvalue and $\beta_t^2 \equiv e^{-2t}\signoise^2 + \Delta_t$
for the null-block eigenvalue, the population $M_t$ has the block-diagonal
spectrum
\(\mathrm{spec}(M_t) = \{\alpha_t^2\}^{\dint} \cup \{\beta_t^2\}^{\dlat-\dint}\)
in the $Q$-basis (the cluster signal subspace and its orthogonal
complement).

\paragraph{Linear part: signal and noise-dim bulks.}
\begin{lemma}[Rank and block structure of $\Umat^{\mathrm{lin}}$]
\label{lem:Ulin}
$\mathrm{rank}(\Umat^{\mathrm{lin}}) = \dlat$ almost surely, and its
nonzero spectrum is the spectrum of $M_t$ rescaled by
$\mu_1^{\,2}/\psi_p$ up to Marchenko--Pastur fluctuations of
relative size $O(1/\sqrt{\psi_p})$. In particular it splits into
a \emph{signal bulk} of size $\dint$ with eigenvalues
$\Theta(\mu_1^{\,2}\alpha_t^{\,2}/\psi_p)$ and a \emph{noise-dim bulk}
of size $\dlat - \dint$ with eigenvalues
$\Theta(\mu_1^{\,2}\beta_t^{\,2}/\psi_p)$.
\end{lemma}
\begin{proof}
$W$ has i.i.d.\ Gaussian entries, so $W^\top W/\dlat$ is full rank a.s.\
for $\pwidth\ge\dlat$, with $\E[W^\top W/\dlat] = \psi_p\,\mathbb{I}_{\dlat}$
and operator-norm fluctuations $O(1/\sqrt{\psi_p})$
\citep{pennington2017nonlinear}. The nonzero eigenvalues of
$WM_tW^\top$ equal those of $M_t^{1/2}W^\top W M_t^{1/2}$, which are
the eigenvalues of $M_t$ rescaled by $\psi_p$ up to those fluctuations.
The block counts and eigenvalue magnitudes follow from \eqref{eq:Mt}.
\end{proof}

\paragraph{Diagonal part: the sample bulk.}
\begin{lemma}[Rank and edge of $\Umat^{\mathrm{diag}}$]
\label{lem:Udiag}
In the proportional limit,
$\mathrm{rank}(\Umat^{\mathrm{diag}}) = \nsamp$ almost surely, and its
nonzero eigenvalues are concentrated around
$\bar\eta_\star/\psi_p$ where
$\bar\eta_\star = \E_{x\sim P_t}[\eta_\star(x,x)]
 = \sum_{k\ge 3}\mu_k^{\,2}\,\E[(\|x\|^2/\dlat)^k]$
with $x \sim P_t$ a sample of the diffused mixture distribution.
\end{lemma}
\begin{proof}
$\Umat^{\mathrm{diag}}$ in \eqref{eq:U-diag} is a sum of $\nsamp$
rank-one outer products $v_\mu v_\mu^\top$ scaled by
$\Lambda_\mu = \eta_\star(x_t^\mu,x_t^\mu)/(\pwidth\nsamp)$,
where $v_\mu = \phi^\perp(x_t^\mu)/\|\phi^\perp(x_t^\mu)\|$.
The kernel decomposition \eqref{eq:kernel-decomp} gives
$v_\mu^\top v_\nu = \delta_{\mu\nu} + o_d(1)$, so the $\{v_\mu\}$
are asymptotically orthonormal and the rank is $\nsamp$ for
$\nsamp < \pwidth$. The eigenvalues are
$\Lambda_\mu \cdot \|\phi^\perp(x_t^\mu)\|^2 \to
\eta_\star(x_t^\mu,x_t^\mu)^2/(\pwidth\nsamp)$ and concentrate
around their mean as a self-averaging quadratic form in the
Gaussian noise component of $x_t^\mu$.
\end{proof}

\begin{remark}[Why a linear-only derivation cannot resolve the sample bulk]
$\Umat^{\mathrm{lin}}$ has rank exactly $\dlat$, so any spectral
content beyond the two data-side blocks must come from a different
contribution. The Hermite-$k\!\ge\!3$ terms in \eqref{eq:kernel-decomp}
are the smallest piece of $\tanh$ on off-diagonal kernel entries but
the \emph{only} piece that survives on the diagonal, and that
survival is what produces the rank-$\nsamp$ sample bulk. An earlier
version of this derivation tried to extract the sample bulk from
the rank-$\nsamp$ fluctuation $\hat S - \Sigmadata$ in
\eqref{eq:Mt}; those modes lie inside the rank-$\dlat$ range of
$\Umat^{\mathrm{lin}}$ and contribute Marchenko--Pastur deformations
of the two data-side bulks, not a separately gapped class.
\end{remark}

\begin{remark}[Fixed-$W$ vs.\ $W$-averaged kernel]
\label{rem:fixedW}
The empirical $\Umat$ averages over the diffusion noise and training
data at \emph{fixed} $W$, so the kernel decomposition
\eqref{eq:kernel-decomp}, which holds in expectation over $W$, must be
upgraded to a concentration statement. By
\citet{benigni2021eigenvalue} for sub-Gaussian entries of $W$ (in
particular Gaussian, our case) the off-diagonal kernel entries
$\phi(x_t^\mu)^\top\phi(x_t^\nu)$ concentrate around their $W$-mean
at rate $O(1/\sqrt{\pwidth})$ in the proportional limit, and the
diagonal entries concentrate at the same rate. The rank-$\nsamp$
structure of $\Umat^{\mathrm{diag}}$ therefore survives at fixed $W$
up to operator-norm corrections of order $\nsamp/\sqrt{\pwidth}=
\sqrt{\psi_n/(\psi_p \pwidth)}\to 0$, which is the self-averaging
argument referenced implicitly in the Lemma 2 proof.
\end{remark}

\begin{remark}[Marchenko--Pastur spreading of the bulk edges]
\label{rem:MP}
Each of the four bulks in Theorem~\ref{thm:fourbulk} is a single
\emph{leading-order} edge; the actual bulk in finite $\dlat$ is a
Marchenko--Pastur-deformed band around that edge with relative width
$\Delta_{\mathrm{MP}} = (\sqrt{\psi_p}+1)^2/(\sqrt{\psi_p}-1)^2 - 1
\approx 4/\sqrt{\psi_p}$ at large $\psi_p$. At our default
$\psi_p = 64$ this gives $\Delta_{\mathrm{MP}}\approx 0.50$, so the
two data-side bulks remain spectrally resolvable whenever
$\alpha_t^{\,2}/\beta_t^{\,2}\gtrsim 1.5$; the empirical
$\alpha_t^{\,2}/\beta_t^{\,2}\ge 4$ at $\signoise=0.5$ and
$\ge 10^4$ at $\signoise=0.01$ satisfy this comfortably. The
noise-dim-to-sample gap is governed instead by the relative size of
the Hermite-1 and Hermite-$\ge 3$ coefficient masses; see the
``Bulk-gap scaling'' paragraph below.
\end{remark}

\paragraph{Four-bulk theorem.}
\begin{theorem}[Four-bulk structure]
\label{thm:fourbulk}
Under the setup above and in the proportional asymptotic limit, the
spectrum of $\Umat$ in \eqref{eq:U-split} splits almost surely into
four contiguous bulks with the following counts and leading-order
edges (fixed $t \in (0,\infty)$):
\begin{center}
\begin{tabular}{@{}lll@{}}
\toprule
Bulk & Count & Leading-order edge \\
\midrule
Signal      & $\dint$                  & $\mu_1^{\,2}\,\alpha_t^{\,2}/\psi_p
                                          \;=\; \mu_1^{\,2}\bigl[e^{-2t}(s^2/\dint+\sigsig^2)+\Delta_t\bigr]/\psi_p$ \\
Noise-dim   & $\dlat-\dint$            & $\mu_1^{\,2}\,\beta_t^{\,2}/\psi_p
                                          \;=\; \mu_1^{\,2}\bigl[e^{-2t}\signoise^{\,2}+\Delta_t\bigr]/\psi_p$ \\
Sample      & $\nsamp$                 & $\bar\eta_\star/\psi_p$ \\
Rank-null   & $\pwidth-\dlat-\nsamp$   & $o(1)$ \\
\bottomrule
\end{tabular}
\end{center}
\end{theorem}
\begin{proof}
Combine Lemmas \ref{lem:Ulin} and \ref{lem:Udiag} with the
decomposition \eqref{eq:U-split}. The ranges of $\Umat^{\mathrm{lin}}$
and $\Umat^{\mathrm{diag}}$ are asymptotically $W$-orthogonal: the
former lies in the column space of $W$ (rank $\dlat$), and the latter
is spanned by the $\nsamp$ residual directions $\phi^\perp(x_t^\mu)$
which are orthogonal to that column space in expectation. Together
they span a subspace of dimension $\dlat + \nsamp < \pwidth$, leaving
a rank-null tail of dimension $\pwidth - \dlat - \nsamp$ with
eigenvalues vanishing as $o_d(1)$ from $R_t$ in \eqref{eq:U-split}.
\end{proof}

\paragraph{Bulk-gap scaling.}
The signal-to-noise-dim gap is the eigenvalue ratio
$\alpha_t^{\,2}/\beta_t^{\,2}
 = \bigl(e^{-2t}(s^2/\dint+\sigsig^2)+\Delta_t\bigr)/\bigl(e^{-2t}\signoise^{\,2}+\Delta_t\bigr)$.
At small diffusion times this is approximately
$(s^2/\dint+\sigsig^2)/\signoise^{\,2}$ (large at $\signoise=0.01$,
modest at $\signoise=0.5$); as $t$ grows both bulks are dragged
toward the diffusion floor $\Delta_t$ and the ratio collapses to $1$.
The noise-dim-to-sample gap is the ratio
$\mu_1^{\,2}\beta_t^{\,2}/\bar\eta_\star$, governed by the relative
sizes of the linear-Hermite coefficient and the higher-Hermite
residual mass: for $\tanh$, $\mu_1^{\,2} \approx 0.367$ while
$\bar\eta_\star$ is dominated by $\mu_3^{\,2} \approx 0.0097$ at our
operating scales, so for $\beta_t^{\,2} \gtrsim \bar\eta_\star/\mu_1^{\,2}$
the noise-dim edge sits above the sample edge and the four bulks
are spectrally resolvable. Below that threshold (small $\signoise$ at
large $t$) the two merge, which is the regime where the four-bulk
cliff at index $\dlat$ becomes a smooth step in
Fig.~\ref{fig:appendix-fourbulk-cliffs}.

\paragraph{Toward a full Stieltjes derivation.}
Theorem \ref{thm:fourbulk} gives leading-order edge magnitudes that
match (within Marchenko--Pastur fluctuations $O(1/\sqrt{\psi_p})$)
the empirical bulk locations in Table~\ref{tab:bulk-sizes} and
Fig.~\ref{fig:bulk-sizes-distances}. A full Stieltjes-transform /
replica derivation that produces the bulk edges as exact roots of a
polynomial fixed-point equation, extending the isotropic-$\Sigmadata$
analysis of \citet{bonnaire2025} to the block-diagonal case using the
tools of \citet{ELKaroui2010spectrum} and
\citet{benigni2021eigenvalue}, is left to follow-up work. The route
taken above is sufficient for the buffer-mechanism corollary in
Section~\ref{sec:rfnn}, which only requires the bulk counts and
leading-order edges established here.


\section{RFNN: eigenvalue saturation and the $\tanh$ pitfall}
\label{app:saturation}

We observed that for sufficiently large data scale $s$, the spectrum of
$\Umat$ collapses to a single bulk with eigenvalues bunched at
$\lambda \approx 1$. The cause is $\tanh$ saturation: when
$\|Wx\|/\sqrt{\dlat} \gtrsim 2$, $\tanh$ outputs $\pm 1$ almost
deterministically and the feature map becomes effectively constant in
$x$. The constraint $s/\sqrt{\dlat} \lesssim 2$ defines the safe
operating regime; all reported sweeps satisfy this constraint.

\section{RFNN: empirical bulk detection}
\label{app:bulk-detection}

Fig.~\ref{fig:bulk-sizes-distances} shows the $\signoise = 0.01$
$\dlat$ sweep with the empirical bulk-detection overlay (peaks
found by \texttt{scipy.signal.find\_peaks} on the histogram counts,
capped at 4 by prominence). At the larger anisotropy ratio
$\sigsig^2/\signoise^2 \approx 10^4$ the four bulks separate
cleanly and the detected peaks land where the theoretical
index-based cuts in Fig.~\ref{fig:fourbulk-hist} put them.

\begin{figure}[!t]
  \centering
  \includegraphics[width=\linewidth]{figures/bulk_detection_hist.pdf}
  \caption{\textbf{Empirical bulk detection across the $\dlat$
    sweep at $\signoise = 0.01$.} Each colored region is one
    detected peak in the eigenvalue density; four bulks are
    detected in every panel. As $\dlat / \dint$ grows, the bulks
    pull apart: at $\dlat = \dint$ the four peaks sit close
    together with shallow valleys, and by $\dlat = 40$
    ($\dlat / \dint = 8$) they have separated into well-isolated
    peaks decades apart in eigenvalue. Per-panel bulk sizes are
    tabulated in Table~\ref{tab:bulk-sizes}; detection algorithm
    in Appendix~\ref{app:methods}.}
  \label{fig:bulk-sizes-distances}
\end{figure}

\section{RFNN: additional 4-bulk spectrum figures}
\label{app:appendix-figures}

The body shows the $\signoise = 0.5$ density view
(Figure~\ref{fig:fourbulk-hist}); for completeness we include the
remaining three corners of the $\{\text{cliffs}, \text{density}\}
\times \{\signoise = 0.01, 0.5\}$ grid here.
Figures~\ref{fig:appendix-fourbulk-cliffs}
and~\ref{fig:appendix-fourbulk-cliffs-sn05} are the cliffs (sorted-spectrum)
view at the two noise scales, and
Figure~\ref{fig:appendix-fourbulk-hist-sn001} is the density view
at $\signoise = 0.01$.
Figure~\ref{fig:appendix-fourbulk-dint} sweeps $\dint$ instead of $\dlat$.

\begin{figure}[!t]
  \centering
  \includegraphics[width=\linewidth]{figures/four_bulk_dlat_sweep.pdf}
  \caption{\textbf{Cliffs view, $\signoise = 0.01$.}
    Sorted log-spectrum across the $\dlat$ sweep. Red and green
    dashed lines mark indices $\dint$ and $\dlat$. The
    $\textsc{bulk\_data\_noise}$ region widens between them as
    $\dlat$ grows.}
  \label{fig:appendix-fourbulk-cliffs}
\end{figure}

\begin{figure}[!t]
  \centering
  \includegraphics[width=\linewidth]{figures/four_bulk_dlat_sweep_sn05.pdf}
  \caption{\textbf{Cliffs view, $\signoise = 0.5$.}
    Same sweep as Figure~\ref{fig:appendix-fourbulk-cliffs} at the
    higher noise scale: the signal/noise-dim cliff at index $\dint$
    is sharp, while the data/sample cliff at index $\dlat$ softens
    as the gap $\sigsig^2/\signoise^2$ shrinks.}
  \label{fig:appendix-fourbulk-cliffs-sn05}
\end{figure}

\begin{figure}[!t]
  \centering
  \includegraphics[width=\linewidth]{figures/four_bulk_dlat_hist_sn001.pdf}
  \caption{\textbf{Density view, $\signoise = 0.01$.}
    Counterpart to the body's Figure~\ref{fig:fourbulk-hist}. The
    larger $\sigsig^2/\signoise^2 \approx 10^4$ pushes the signal
    bulk decades to the right of the noise-dim bulk.}
  \label{fig:appendix-fourbulk-hist-sn001}
\end{figure}

\begin{figure}[!t]
  \centering
  \includegraphics[width=\linewidth]{figures/four_bulk_dint_sweep.pdf}
  \caption{\textbf{Varying $\dint$} at $\dlat = 20$,
    $\signoise = 0.5$. As $\dint \to \dlat$ the noise-dim bulk
    shrinks toward zero width and we recover Bonnaire's two-bulk
    picture.}
  \label{fig:appendix-fourbulk-dint}
\end{figure}
